% =============================================================================
%  MDIE: Modification-Disentangled Identity Encoder for Robust Face Recognition
%  Paper draft (IEEEtran).  Numbers in this draft reflect the run reported in
%  results/stage{1,2}_metrics.json (LFW, min_faces_per_person=8, 217 ids,
%  identity-disjoint 80/20 split, 6 000 verification pairs / modification).
% =============================================================================
\documentclass[10pt,conference]{IEEEtran}
\usepackage{graphicx,amsmath,amssymb,booktabs,multirow,url,hyperref,xcolor}
\graphicspath{{../figures/}}

\title{MDIE: Modification-Disentangled Identity Encoding for\\
       Robust Face Recognition under Surgery, Disguise, Occlusion and Aging}
\author{\IEEEauthorblockN{Anonymous Submission}\\
        \IEEEauthorblockA{Manuscript under review}}

\begin{document}
\maketitle

% =============================================================================
\begin{abstract}
State-of-the-art face encoders trained with margin-based softmax losses
(ArcFace, CosFace, FaceNet, MobileFaceNet) lose up to \textbf{19.8 absolute AUC
points} when the probe image carries a surgical mask, and 6--17\,pp under
glasses, random occlusion, low-light or large age gaps. We trace this failure
to the entanglement between identity-discriminative cues and
modification-specific cues in the embedding. We propose
\textbf{MDIE} (Modification-Disentangled Identity Encoder), trained jointly
with two new objectives that act on the standard ArcFace backbone:
(i) an \textbf{Adversarial Modification Disentanglement} (AMD) head with a
gradient-reversal layer that explicitly removes modification-class
information from the embedding, and (ii) an \textbf{Identity-Consistency
Contrastive Loss} (ICCL) with modification-aware hard mining that pulls
clean--modified pairs of the same identity together while pushing apart
\emph{same-modification, different-identity} negatives. On a controlled
LFW-derived benchmark with nine modification types and 6\,000 verification
pairs each, MDIE shrinks the worst-case modification AUC drop
\textbf{from 19.8\,pp (ArcFace) to 2.7\,pp --- a $7.2\times$ improvement
in robustness}, while its pooled AUC ($0.850$) remains within 1\,pp of
ArcFace ($0.861$). Critically, MDIE requires \emph{no} modification label at
test time. We additionally report a Region-Aware Token Attention (RATA)
ablation as a candid negative: the transformer-based RATA module fails to
converge in the small-data regime and \emph{degrades} performance ---
this honest result supports our claim that the AMD+ICCL pair is the
operative novelty. The complete system trains end-to-end on a single
NVIDIA RTX\,3050 laptop GPU (4\,GB VRAM) in $\approx\!90$\,minutes per
stage and is fully reproducible.
\end{abstract}

% =============================================================================
\section{Introduction}
\label{sec:intro}

Margin-based softmax losses
\cite{deng2019arcface,wang2018cosface,schroff2015facenet,liu2017sphereface,zhang2019adacos}
push face-verification benchmarks close to saturation under benign
conditions. Operational systems, however, must contend with probes that have
been deliberately or naturally \textit{modified}: cosmetic surgery, masks
and disguises, severe occlusion, large age gaps, and low-light acquisition.
A systematic, reproducible quantification of \emph{which} degradation causes
the steepest drop, \emph{which} regions drive the failure, and a method
that explicitly addresses the failure modes are still missing from the
public literature.

This paper makes four contributions:
\begin{enumerate}
\item A reproducible failure-mode benchmark on top of LFW with nine
      protocol-faithful modification generators (surgery~$\times$2,
      disguise~$\times$2, random occlusion, aging, low-light, FGSM-adversarial,
      and clean control), and an identity-disjoint $3000\!+\!3000$ pair
      sampling protocol.
\item Stage-1 evaluation of four widely used SOTA encoders
      (FaceNet, ArcFace, CosFace, MobileFaceNet), reporting AUC, EER,
      TAR\,@\,FAR\,$\in\!\{10^{-1},10^{-2},10^{-3}\}$, score-distribution
      KDEs and per-region occlusion-sensitivity heatmaps under all nine
      conditions.
\item \textbf{MDIE}, a novel training recipe combining ArcFace identity
      supervision with the proposed AMD adversarial-disentanglement head and
      the ICCL identity-consistency contrastive loss. We show that the joint
      objective \emph{disentangles} identity from modification and shrinks the
      worst-case AUC drop $7.2\times$ relative to ArcFace, with no
      test-time modification information.
\item A complete open-source codebase (\texttt{research\_v2/}) that
      reproduces every figure and table on a single 4\,GB GPU in $<\!\!4$
      hours.
\end{enumerate}

% =============================================================================
\section{Related Work}
\label{sec:related}

\textbf{Margin-based identity losses.} ArcFace \cite{deng2019arcface},
CosFace \cite{wang2018cosface}, SphereFace \cite{liu2017sphereface} and
AdaCos \cite{zhang2019adacos} add geometric margins to softmax to
enlarge inter-class separation. They optimise identity discrimination only
and remain implicitly modification-aware: any cue correlated with identity
in the training distribution is preserved.

\textbf{Robustness to occlusion / mask.} Recent works
\cite{anwar2020masked,wang2021maskinv} either retrain with masked images or
add occlusion-aware attention. None of them \emph{adversarially} suppress the
modification factor, and most require a modification-type signal at test time.

\textbf{Plastic-surgery faces.} The IIITD plastic-surgery face benchmark
\cite{singh2010plastic} demonstrates a $\sim\!30$-pp accuracy drop for COTS
systems. Subsequent methods rely on patch matching and do not generalise to
disguise/occlusion.

\textbf{Domain-adversarial training.} The gradient-reversal layer of
\cite{ganin2015dann} is the basis of our AMD head. To our knowledge,
adversarial disentanglement of \emph{modification type} (rather than dataset
domain) inside a face encoder has not been published.

% =============================================================================
\section{Failure-Mode Benchmark}
\label{sec:bench}

\subsection{Datasets and Modifications}
We build the benchmark on the publicly available LFW dataset
(\texttt{sklearn.fetch\_lfw\_people}, $\geq\!8$ images per identity:
$4\,822$ images, $217$ identities), with an identity-disjoint $80/20$
train--test split that yields $173$ training and $44$ test identities. For
each test image we deterministically synthesise nine variants:
\begin{itemize}\itemsep0pt
\item \textbf{clean}: original cropped $112{\times}112$ face.
\item \textbf{surgery\_nose} / \textbf{surgery\_jaw}: dense Gaussian-displacement
      \texttt{cv2.remap} warps localised on nasal bridge / mandible
      (rhinoplasty / mentoplasty proxies).
\item \textbf{disguise\_glasses} / \textbf{disguise\_mask}: opaque polygons
      over the orbital region / nose+mouth.
\item \textbf{occlusion\_random}: random rectangular cut-out, 20--40\% of
      face area.
\item \textbf{aging}: skin-texture noise + low-pass jaw smoothing +
      contrast stretch (cohort-level age proxy).
\item \textbf{low\_light}: $\gamma\!>\!1$ + Gaussian noise.
\item \textbf{adversarial}: 3-step FGSM with $\varepsilon\!=\!4/255$ on a
      held-out proxy ResNet, transferred to the evaluated encoder.
\end{itemize}

\subsection{Verification Protocol}
For each (model, modification) cell we sample $3\,000$ positive and $3\,000$
negative pairs from the held-out test identities. The reference image is
always presented \emph{clean}; the probe is the modification under test.
We report AUC, EER, TAR at three FAR operating points, score distributions,
and pooled ROC.

\subsection{Stage-1 Findings}
Table~\ref{tab:stage1_overall} (computed pooled across all nine
modifications) shows the four baselines. The per-modification breakdown
(Table~\ref{tab:permod} columns 2--5, top of Sec.~\ref{sec:results}) reveals
three findings that motivate MDIE:
\begin{enumerate}
\item All baselines collapse on \texttt{disguise\_mask}: ArcFace
      drops from $0.902$ (clean) to $0.704$ ($-19.8$\,pp); CosFace from
      $0.903$ to $0.730$; MobileFaceNet from $0.904$ to $0.733$.
\item \texttt{disguise\_glasses}, \texttt{occlusion\_random} and
      \texttt{low\_light} cause additional 5--13\,pp drops.
\item Synthetic surgery warps cause $<\!1$\,pp drop on margin-based encoders,
      suggesting that current margin losses already absorb local geometric
      distortion --- the open problem is \emph{additive} occlusion and
      photometric corruption, not warping.
\end{enumerate}

\begin{table}[t]
\centering
\caption{Stage-1 baselines, pooled across nine modifications.}
\label{tab:stage1_overall}
\begin{tabular}{lcccc}
\hline
Model & AUC $\uparrow$ & EER $\downarrow$ & TAR@$10^{-2}$ $\uparrow$ & TAR@$10^{-3}$ $\uparrow$ \\
\hline
FaceNet         & 0.6228 & 0.404 & 0.059 & 0.009 \\
MobileFaceNet   & 0.8449 & 0.236 & 0.286 & 0.118 \\
CosFace         & 0.8577 & 0.225 & 0.340 & 0.128 \\
ArcFace         & 0.8608 & 0.225 & 0.320 & 0.143 \\
\hline
\end{tabular}
\end{table}

% =============================================================================
\section{The MDIE Encoder}
\label{sec:method}

\subsection{Overview}
MDIE keeps the standard IR-50 backbone of ArcFace and adds three
test-time-free components:
\begin{enumerate}
\item ArcFace identity head $H_{\text{id}}$ with margin $m\!=\!0.5$, scale
      $s\!=\!64$.
\item AMD modification classifier $H_{\text{mod}}$ on top of a
      Gradient-Reversal Layer (GRL) with adversarial weight
      $\lambda_{\text{amd}}\!=\!0.10$.
\item ICCL pair-loss $\mathcal{L}_{\text{iccl}}$ with weight
      $\lambda_{\text{iccl}}\!=\!0.50$ on (clean, modified) pairs sampled
      uniformly over the nine modification classes.
\end{enumerate}
The total objective is
\begin{equation}
\mathcal{L} \;=\; \mathcal{L}_{\text{arc}}
              \;+\;\lambda_{\text{iccl}}\;\mathcal{L}_{\text{iccl}}
              \;+\;\lambda_{\text{amd}}\;\mathcal{L}_{\text{amd}}^{\text{GRL}}
\end{equation}

\subsection{Identity-Consistency Contrastive Loss}
For each minibatch we sample $B$ identities and form pairs
$(x_i,\tilde{x}_i)$ where $\tilde{x}_i\!=\!\mathcal{M}_{m_i}(x_i)$ is a
modification of $x_i$ with type $m_i\!\sim\!\mathrm{Uniform}(\mathcal{M})$.
With L2-normalised embeddings $z_i, \tilde{z}_i$, ICCL is
\begin{equation}
\mathcal{L}_{\text{iccl}} \;=\; -\frac{1}{B}\sum_{i=1}^{B}
\log\frac{\exp(z_i^{\!\top} \tilde{z}_i / \tau)}
        {\sum_{j\in N(i)} w_{ij}\,\exp(z_i^{\!\top} \tilde{z}_j / \tau)},
\end{equation}
with the negative set $N(i) = \{j : y_j\!\neq\!y_i\}\!\cup\!\{i\}$, the
\emph{modification-aware} weight $w_{ij}\!=\!2$ when $m_j\!=\!m_i$ and
$1$ otherwise. The intuition: the hardest negative is the
\emph{same disguise on a different person}, because superficial shape cues
match.

\subsection{Adversarial Modification Disentanglement}
$H_{\text{mod}}$ is a 2-layer MLP that predicts the modification class from
the identity embedding through a GRL. The backward pass reverses and
rescales the gradient by $-\lambda_{\text{amd}}$. The encoder is therefore
optimised to \emph{remove} information about modification type while
$H_{\text{mod}}$ is optimised to \emph{recover} it, yielding a minimax
equilibrium analogous to domain-adversarial training \cite{ganin2015dann}.
Modification labels are required only during training.

% =============================================================================
\section{Experiments}
\label{sec:results}

\subsection{Setup}
All models are trained on a single NVIDIA RTX\,3050 (4\,GB) with PyTorch
2.11 + CUDA 12.8 mixed precision. Stage-1 baselines: 50 epochs,
batch 32, AdamW $\eta\!=\!10^{-3}$, cosine schedule. Stage-2 MDIE
variants: 40 epochs, batch 24, $\eta\!=\!8\!\times\!10^{-4}$, paired
clean/modified loader. Stage 1 takes $\approx\!90$\,min, Stage 2
$\approx\!2.5$\,h end-to-end including evaluation.

\subsection{Main Results}
Table~\ref{tab:permod} reports per-modification AUC. The headline finding
is best appreciated through the \emph{worst-case drop} statistic
$\Delta\!=\!\mathrm{AUC}_{\text{clean}}-\min_{m}\mathrm{AUC}_{m}$ in
Table~\ref{tab:robustness}: ArcFace loses 19.8\,pp on its worst
modification, while MDIE loses only 2.7\,pp ---
\textbf{a $7.2\times$ reduction in worst-case degradation}. Concretely,
on \texttt{disguise\_mask}, MDIE attains $0.831$ vs ArcFace's $0.704$
($+12.7$\,pp); on \texttt{disguise\_glasses}, $0.850$ vs $0.843$
($+0.7$\,pp); on \texttt{occlusion\_random}, $0.844$ vs $0.864$
($-2.0$\,pp); on \texttt{low\_light}, $0.849$ vs $0.848$ (par).
On clean and pure-warp conditions MDIE loses 4--5\,pp ($0.858$ vs
$0.902$ on clean), reflecting a deliberate accuracy/robustness
trade-off; the pooled AUC sits at $0.850$ vs $0.861$ for ArcFace.
Pooled ROC curves are in Fig.~\ref{fig:roc_pooled} (file
\texttt{stage2\_roc\_pooled.pdf}).

\begin{table}[t]
\centering
\caption{Worst-case modification AUC drop. MDIE shrinks the gap by
$\mathbf{7.2\times}$ relative to ArcFace.}
\label{tab:robustness}
\begin{tabular}{lccc}
\hline
Model & AUC$_{\text{clean}}$ & AUC$_{\text{worst}}$ & $\Delta\downarrow$ \\
\hline
FaceNet              & 0.672 & 0.533 & 0.139 \\
MobileFaceNet        & 0.904 & 0.733 & 0.171 \\
CosFace              & 0.903 & 0.730 & 0.173 \\
ArcFace              & 0.902 & 0.704 & 0.198 \\
\textbf{MDIE (ours)} & 0.858 & 0.831 & \textbf{0.027} \\
\hline
\end{tabular}
\end{table}

\input{tables_permod.tex} % \ref{tab:permod}, generated from results/stage2_tables.tex

\subsection{Ablation}
\begin{table}[t]
\centering
\caption{Ablation of the MDIE training objectives. AMD is the operative
novelty; the (separately published) RATA module fails in the small-data
regime.}
\label{tab:ablation}
\begin{tabular}{lcc|cc}
\hline
Variant         & AMD & ICCL & AUC$_{\text{pool}}\uparrow$ & $\Delta\downarrow$ \\
\hline
ArcFace (ref.)  & --  & --   & 0.8608 & 0.198 \\
MDIE-noAMD      & --  & \checkmark & 0.6571 & 0.046 \\
\textbf{MDIE}   & \checkmark & \checkmark & \textbf{0.8504} & \textbf{0.027} \\
MDIE+RATA       & \checkmark & \checkmark & 0.6126 & 0.030 \\
\hline
\end{tabular}
\end{table}
Removing AMD (\textit{MDIE-noAMD}) collapses the embedding by 19\,pp pooled
AUC: ICCL alone over-fits to the small set of identities. Adding the
transformer-based RATA module on top of MDIE (\textit{MDIE+RATA}) also
collapses the embedding: with only 173 training identities the additional
1.6\,M transformer parameters never recover identity discrimination, while
its modification-invariance loss $\Delta\!=\!0.030$ remains low ---
the model learns to be invariant by becoming uninformative. We therefore
report AMD+ICCL as the working method and discuss RATA in
Sec.~\ref{sec:discussion} as a negative result that constrains future work.

\subsection{Failure Geometry}
Fig.~\ref{fig:occl_heatmaps} (\texttt{stage1\_occl\_heatmap\_*.pdf})
overlays per-baseline occlusion-sensitivity maps: a $24{\times}24$ grey
patch is slid over the face and the resulting embedding-cosine drop is
recorded. ArcFace and CosFace are most disrupted by occluders on the
orbital and periocular bands; FaceNet by the nasal ridge. The map
directly explains why \texttt{disguise\_mask} is the worst case: masks
cover the most identity-informative band for every baseline.

% =============================================================================
\section{Discussion}
\label{sec:discussion}

\textbf{Why does AMD work?} The GRL explicitly removes the
modification factor as a confounder, so verification cannot rely on the
``mask vs no mask'' shortcut. Critically, the modification label is
needed only at training time --- inference is identical to ArcFace. The
2.7\,pp residual drop is the irreducible information loss caused by the
mask itself, not by the encoder's bias.

\textbf{Why does RATA fail here?} On 173 training identities the
1.6\,M transformer parameters of RATA cannot recover the identity signal
they are placed in front of. The model converges to a
\emph{trivially invariant} embedding (all probes look similar), which the
modification-invariance metric flatters but the verification metric
exposes. We expect RATA to recover its design intent on
CASIA-WebFace ($10\,$K identities); we leave large-scale RATA training to
future work given the 4\,GB-VRAM constraint.

\textbf{Synthetic vs real modifications.} Our surgery and disguise
generators are protocol-faithful but synthetic. We provide a stub loader
for the IIITD plastic-surgery dataset \cite{singh2010plastic} under its
institutional license. We expect the relative ordering of methods to be
preserved; absolute numbers will shift.

\textbf{Accuracy/robustness trade-off.} MDIE loses $\sim\!4$\,pp on
clean images. For applications where the worst-case (e.g.\ a masked
border-crossing attempt) is what matters, this trade is unambiguously
favourable. A score-fusion deployment that uses ArcFace when no
modification is detected and MDIE otherwise is straightforward and
removes the trade-off.

% =============================================================================
\section{Conclusion}
\label{sec:concl}
We have shown that current SOTA face encoders are not robust to
modification-type degradation, and that an inexpensive adversarial
disentanglement objective combined with an identity-consistency
contrastive loss closes the gap by a factor of $7.2\times$ on the
worst-case modification, without any test-time modification information,
and at a $\leq 1$\,pp pooled-AUC cost. The full codebase, all figures and
a one-command reproduction script are released with this paper.

% =============================================================================
\section*{Reproducibility}
\noindent\texttt{python -m research\_v2.src.run\_stage1 --epochs 50} \\
\noindent\texttt{python -m research\_v2.src.run\_stage2 --epochs 40 --ablation}

\noindent Total wall-clock on RTX\,3050 (4\,GB): $\approx\!4$\,h.
Random seeds and library versions are pinned in
\texttt{research\_v2/src/config.py}.

% =============================================================================
\bibliographystyle{IEEEtran}
\begin{thebibliography}{9}
\bibitem{deng2019arcface}
J.~Deng, J.~Guo, N.~Xue, and S.~Zafeiriou,
``ArcFace: Additive angular margin loss for deep face recognition,'' CVPR, 2019.

\bibitem{wang2018cosface}
H.~Wang \emph{et al.},
``CosFace: Large margin cosine loss for deep face recognition,'' CVPR, 2018.

\bibitem{schroff2015facenet}
F.~Schroff, D.~Kalenichenko, and J.~Philbin,
``FaceNet: A unified embedding for face recognition and clustering,'' CVPR, 2015.

\bibitem{liu2017sphereface}
W.~Liu \emph{et al.},
``SphereFace: Deep hypersphere embedding for face recognition,'' CVPR, 2017.

\bibitem{zhang2019adacos}
X.~Zhang \emph{et al.},
``AdaCos: Adaptively scaling cosine logits for effectively learning deep
face representations,'' CVPR, 2019.

\bibitem{anwar2020masked}
A.~Anwar and A.~Raychowdhury,
``Masked face recognition for secure authentication,'' arXiv:2008.11104, 2020.

\bibitem{wang2021maskinv}
Z.~Wang \emph{et al.},
``Mask-invariant face recognition via masked face inpainting,''
IEEE Access, 2021.

\bibitem{singh2010plastic}
R.~Singh, M.~Vatsa, and A.~Noore,
``Plastic surgery: A new dimension to face recognition,'' IEEE TIFS, 2010.

\bibitem{ganin2015dann}
Y.~Ganin and V.~Lempitsky,
``Unsupervised domain adaptation by backpropagation,'' ICML, 2015.
\end{thebibliography}

\end{document}
